\section{The dataset}
\label{sec:data}

\subsection{Participants, protocol and sensors}

Ten young adults (6 women, 4 men, 20--31 years) walked along a walkway and crossed six obstacles: three heights
(H1, H2, H3 = 10, 20 and 30\,\% of their own leg length) $\times$ two widths (W1 = 5\,cm, W2 = 15\,cm deep along the
walkway), about five times each, plus level walks with no obstacle. The obstacles were run in blocks (about five
consecutive repeats of the same obstacle, block order different between participants). Table~\ref{tab:participants}
gives the participants, their nominal obstacle heights in centimetres and the number of rows used.

\begin{itemize}
  \item \textbf{Wearables:} six Movella Dot IMUs (thigh, shank and foot of both legs), sampled at 60\,Hz. They provide the
  segment orientation (Euler angles), acceleration, angular velocity and quaternion.
  \item \textbf{Reference:} a camera motion-capture system (toe and heel markers of both feet, obstacle markers),
  synchronised to the Dot clock with three left-leg raises at the start of the session (step~5 of the pipeline).
  \item \textbf{Leg length} (participant sheet) defines the obstacle heights and normalises lengths.
\end{itemize}

\begin{table}[htbp]
\centering
\caption{Participants and rows of the dataset. H1/H2/H3 = nominal obstacle heights (10/20/30\,\% of the leg length).
``Used (ML)'' = rows with \vn{include_ml} = 1 (Section~\ref{sec:quality}).}
\label{tab:participants}
\begin{adjustbox}{max width=\textwidth}
\small
\begin{tabular}{@{}l c r r r r r r r r r@{}}
\toprule
\textbf{Participant} & \textbf{Sex} & \textbf{Age} & \textbf{Height} & \textbf{Leg} & \textbf{Mass} & \textbf{H1 / H2 / H3} & \textbf{Obstacle} & \textbf{Used} & \textbf{Level} & \textbf{Used} \\
 & & \textbf{(y)} & \textbf{(cm)} & \textbf{length (cm)} & \textbf{(kg)} & \textbf{(cm)} & \textbf{rows} & \textbf{(ML)} & \textbf{rows} & \textbf{(ML)} \\
\midrule
101 & Female & 22 & 164.0 & 90.0 & 94.2 & 9.0 / 18.0 / 27.0 & 32 & 18 & 8 & 5 \\
102 & Male & 23 & 189.0 & 99.0 & 75.8 & 9.9 / 19.8 / 29.7 & 30 & 30 & 5 & 4 \\
103 & Male & 25 & 186.0 & 94.0 & 84.9 & 9.4 / 18.8 / 28.2 & 30 & 29 & 5 & 4 \\
104 & Female & 31 & 162.0 & 86.5 & 65.9 & 8.7 / 17.3 / 25.9 & 30 & 26 & 3 & 3 \\
105 & Female & 20 & 169.0 & 87.5 & 87.7 & 8.8 / 17.5 / 26.2 & 30 & 25 & 4 & 4 \\
106 & Male & 21 & 194.0 & 106.0 & 83.3 & 10.6 / 21.2 / 31.8 & 30 & 29 & 3 & 3 \\
107 & Male & 24 & 191.0 & 95.5 & 98.8 & 9.6 / 19.1 / 28.6 & 30 & 27 & 4 & 4 \\
108 & Female & 20 & 174.0 & 91.5 & 61.7 & 9.2 / 18.3 / 27.4 & 30 & 26 & 4 & 4 \\
109 & Female & 23 & 179.5 & 93.5 & 56.0 & 9.3 / 18.7 / 28.1 & 31 & 26 & 5 & 4 \\
110 & Female & 31 & 168.5 & 83.0 & 101.5 & 8.3 / 16.6 / 24.9 & 30 & 24 & 5 & 4 \\
\midrule
\textbf{All} & & & & & & & 303 & 260 & 46 & 39 \\
\bottomrule
\end{tabular}

\end{adjustbox}
\end{table}

\subsection{From recordings to one row per crossing}
\label{sec:rows}

The dataset is built by \vn{Analysis/ML/build_crossing_master.m}. It only reads the outputs of the processing pipeline
(steps 1--9) and writes \vn{Results/ML Dataset/}. One row is one obstacle crossing (participant $\times$ obstacle $\times$
repeat) or one level walk; both limbs are in the same row, named by their role (\vn{lead_}, \vn{trail_}).

\begin{description}[leftmargin=0.4cm,style=nextline]
  \item[Stride] Zero-velocity point (ZVP) to the next ZVP of the same foot (mid-stance to mid-stance), from the Dot foot
  sensor (step 4).
  \item[Swing] The step-4 toe-off (minimum of the foot sagittal angle after the stride's first ZVP) to the step-4
  heel-strike (maximum of the foot angle) inside that stride. All ``\% swing'' values use this window
  (0 = toe-off, 100 = heel-strike).
  \item[Crossing stride and roles] From the camera crossings file: the leading limb is the foot that passes the obstacle
  centre ($y = 0$) first; for each limb, the crossing stride is the ZVP stride that contains the moment its toe passes
  $y = 0$. The Dot alone recovers the same leading side and strides (\vn{role_imu}, Section~\ref{sec:res-device}).
  \item[Level walks] The strides in which each toe crosses the obstacle line ($y = 0$) during the level passes before the
  first obstacle; the foot that crosses first is called ``lead''.
  \item[Segment angles] Sagittal tilt of the thigh, shank and foot (Dot Euler angle about the medio-lateral axis),
  sign-normalised per participant and side with the step-9 templates: maximum = thigh flexion / shank forward /
  toe-up; minimum = thigh backward / shank backward (heel-up) / toe-down.
  \item[Foot trajectory (ZHC)] The foot sensor's 3-D position, integrated from the acceleration over each stride with a
  zero-velocity update and a level-ground correction (step 4). Height is relative to the stride's first ZVP
  (foot flat), progression along the stride direction.
  \item[Foot acceleration] Magnitude of the foot acceleration in the global frame, gravity removed.
\end{description}

\subsection{Variables}
\label{sec:variables}

The master file has 193 columns (Table~\ref{tab:vargroups}); every column is defined in Appendix~\ref{app:variables}
(Table~\ref{tab:variables}). Naming rules of the wearable features (per limb, prefix \vn{lead_} or \vn{trail_}):

\begin{itemize}
  \item \vn{<segment>_max}, \vn{_min}, \vn{_rom}: largest and smallest tilt in the swing and their difference (deg);
  \vn{_t_max}, \vn{_t_min}: when they occur (\% swing); \vn{_max_corr}, \vn{_min_corr}: the same extremes minus the
  stride's own starting angle (removes the sensor mounting offset).
  \item \vn{<segment>_av_max}, \vn{_av_min}, \vn{_av_t_max}, \vn{_av_t_min}, \vn{_av_toeoff}: angular velocity (deg/s, time
  derivative of the tilt) extremes, their timing and the value at toe-off.
  \item \vn{foot_acc_peak_start}, \vn{foot_acc_peak_land}, \vn{foot_acc_rms_swing}: foot acceleration in the first and last
  25\,\% of the swing and over the whole swing (m/s$^2$).
  \item \vn{zhc_peak_foot_h}: highest foot-sensor height in the swing (mm), with its time \vn{zhc_peak_h_t} and progression
  \vn{zhc_peak_h_x}; \vn{zhc_h_25}, \vn{_50}, \vn{_75}: height at 25/50/75\,\% of the swing; \vn{zhc_h_toeoff},
  \vn{zhc_h_landing}; peak vertical and forward speeds.
  \item Gait: \vn{swing_time}, \vn{stance_time}, \vn{stride_len_app1/2} (the two strides before the crossing),
  \vn{speed_approach}, \vn{speed_cross}, \vn{cadence}.
  \item Camera reference (never an input): minimum clearance, take-off and landing distances, toe-marker peak, obstacle
  height, and the camera toe lift-off time \vn{toeoff_t_cam} (Section~\ref{sec:toeoff}).
\end{itemize}

\begin{table}[htbp]
\centering
\caption{Column groups of \vn{Crossing_Master_YoungAdults.csv} and their roles.}
\label{tab:vargroups}
\small
\begin{tabular}{@{}p{6.2cm} r p{8.2cm}@{}}
\toprule
\textbf{Group} & \textbf{Columns} & \textbf{Roles} \\
\midrule
Identifiers and obstacle labels & 9 & LABEL 7, ID 2 \\
Design and trial order & 7 & META 7 \\
Participant & 6 & META 6 \\
Events, timing and durations & 19 & META 12, FEATURE 7 \\
Segment tilt angles over the swing (Dot) & 42 & FEATURE 42 \\
Segment angular velocity over the swing (Dot) & 30 & FEATURE 30 \\
Foot acceleration (Dot) & 6 & FEATURE 6 \\
Foot trajectory (ZHC) and gait (Dot) & 29 & FEATURE 29 \\
Camera reference values (never model inputs) & 20 & REFERENCE/TARGET 12, REFERENCE 8 \\
Quality flags and inclusion & 25 & FLAG 22, REFERENCE 2, META 1 \\
\midrule
\textbf{All columns} & 193 & FEATURE 114, META 26, FLAG 22, REFERENCE/TARGET 12, REFERENCE 10, LABEL 7, ID 2 \\
\bottomrule
\end{tabular}

\end{table}

\paragraph{Model inputs.} Only FEATURE columns (wearable) are model inputs: 114 in the dictionary, of which 110 are used
(the two step lengths were never computable, and the two toe-height columns are excluded because their sensor-to-toe vector
was fitted with the camera). They split into 57 leading-limb features (including the approach strides), 52 trailing-limb
features and the crossing duration. Additional inputs are derived from the curve files:

\begin{table}[htbp]
\centering
\caption{Feature sets used by the models.}
\label{tab:featuresets}
\small
\begin{tabular}{@{}p{4.4cm} r p{9.4cm}@{}}
\toprule
\textbf{Set} & \textbf{Size} & \textbf{Content} \\
\midrule
Summary features & 110 & All FEATURE columns (57 leading, 52 trailing, 1 shared). \\
Foot-IMU device set & 46 / 41 & Leading / trailing foot: foot tilt and angular velocity, foot acceleration, foot trajectory,
swing and stance time (+ approach strides for the leading foot) and the 13 clearance-shape features. \\
Thigh or shank set & 12 per limb & Tilt and angular-velocity features of that segment. \\
Clearance-shape features & 13 per limb & From the swing curves of foot height and progression: peak height, progression to
the peak, horizontal distance and \% of swing spent above 50/70/80/90\,\% of the peak, foot tilt at the apex, progression
after the apex, rise/fall steepness ratio. \\
Swing curves & $6 \times 51$ per limb & Thigh, shank and foot tilt, foot acceleration, foot height and progression at
0, 2, \ldots, 100\,\% of the swing (\vn{Crossing_Curves_YoungAdults.csv}). \\
Time-aligned signals & $5 \times 127$ per limb & Thigh, shank and foot tilt, foot acceleration and foot height at fixed
samples from $-1.5$\,s to $+0.6$\,s around the trailing toe lift-off (\vn{Crossing_Aligned_YoungAdults.csv}). \\
Stride features & 12 + 6 & For every stride of the session (camera-free pipeline): 12 features of the foot trajectory, and
6 ``online'' features known at the foot apex. \\
\bottomrule
\end{tabular}
\end{table}

\subsection{Quality control and inclusion}
\label{sec:quality}

\paragraph{Inclusion flags.} \vn{include_biomech} = gait events present, Dot data present and no data-quality flag.
\vn{include_ml} additionally requires a plausible foot trajectory (stride length 200--2500\,mm, peak height 10--800\,mm) and no
integration-drift flag. Of 303 obstacle rows, 260 are used (Table~\ref{tab:exclusions}); of 46 level walks, 39.
The 260 crossings are balanced across heights (H1 88, H2 90, H3 82) and widths (W1 127, W2 133).

\begin{table}[htbp]
\centering
\caption{Obstacle rows not used by the ML analyses (first reason that applies).}
\label{tab:exclusions}
\small
\begin{tabular}{@{}p{9.5cm} r@{}}
\toprule
\textbf{First reason that applies} & \textbf{Obstacle rows} \\
\midrule
data-quality flag (stride role or window check failed) & 18 \\
no toe-off / heel-strike inside the crossing stride & 11 \\
implausible foot trajectory (ZHC stride length or peak) & 10 \\
foot-trajectory integration drift & 4 \\
\midrule
\textbf{Excluded from the ML analyses} & \textbf{43} \\
\textbf{Used (\texttt{include\_ml} = 1)} & \textbf{260} \\
\bottomrule
\end{tabular}

\end{table}

\paragraph{Sampling and completeness.} The Dot recordings were audited for this report (Table~\ref{tab:sampling}): the
time base is exactly 60\,Hz in every test, with no gaps or repeated samples, no dropped packets and no missing Dot samples
inside any crossing.

\begin{table}[htbp]
\centering
\caption{Dot sampling audit (build report, task 14).}
\label{tab:sampling}
\begin{adjustbox}{max width=\textwidth}
{\footnotesize
\begin{tabular}{lrrrrrrrrr}
\toprule
\textbf{participant} & \textbf{fs (Hz)} & \textbf{samples} & \textbf{sampling interval min-max (ms)} & \textbf{time gaps} & \textbf{Dot sensors} & \textbf{dropped / missing packets} & \textbf{NaN angle samples (worst sensor)} & \textbf{NaN acc samples (worst sensor)} & \textbf{crossings with a NaN Dot angle} \\
\midrule
101 & 60 & 39114 & 16.67--16.67 & 0 & 6 & 0 & 0.000\% & 0.000\% & 0 \\
102 & 60 & 37805 & 16.67--16.67 & 0 & 6 & 0 & 0.000\% & 0.000\% & 0 \\
103 & 60 & 37572 & 16.67--16.67 & 0 & 6 & 0 & 0.000\% & 0.000\% & 0 \\
104 & 60 & 47715 & 16.67--16.67 & 0 & 6 & 0 & 0.000\% & 0.000\% & 0 \\
105 & 60 & 58252 & 16.67--16.67 & 0 & 6 & 0 & 0.000\% & 0.000\% & 0 \\
106 & 60 & 44381 & 16.67--16.67 & 0 & 6 & 0 & 0.000\% & 0.000\% & 0 \\
107 & 60 & 41291 & 16.67--16.67 & 0 & 6 & 0 & 0.000\% & 0.000\% & 0 \\
108 & 60 & 39019 & 16.67--16.67 & 0 & 6 & 0 & 0.000\% & 0.000\% & 0 \\
109 & 60 & 38737 & 16.67--16.67 & 0 & 6 & 0 & 0.000\% & 0.000\% & 0 \\
110 & 60 & 44144 & 16.67--16.67 & 0 & 6 & 0 & 0.000\% & 0.000\% & 0 \\
\bottomrule
\end{tabular}
}
\end{adjustbox}

\end{table}

\paragraph{Test 106 thigh values.} In Test 106 a few thigh values were outliers (more than $3 \times 1.4826$ median
absolute deviations from the participant $\times$ obstacle $\times$ role median, the step-9 rule). They are set to empty in the
dataset (26 values in 13 crossings, all logged with the original value in \vn{exclusions_log.csv}), and the thigh curve of
the same limb is left empty in the curve files.

\paragraph{Toe-off timing.}
\label{sec:toeoff}
The step-4 toe-off is the minimum of the foot angle. On the trailing limb of an obstacle crossing the foot keeps pointing
its toe down while it lifts, so this minimum comes well after the toe has left the floor (Table~\ref{tab:toeoff}): the
camera toe marker leaves the floor a median 233\,ms earlier, when the foot sensor is already 356\,mm up at the IMU event.
The IMU events are kept (they define all swing features consistently with steps 4--9), and a camera toe lift-off
(\vn{lead_toeoff_t_cam}, \vn{trail_toeoff_t_cam}: first sample at which the toe marker is more than 10\,mm above its stance
height for four consecutive samples) was added for timing questions. In this report the deadline of a cue for the
trailing limb is always the camera toe lift-off.

\begin{table}[htbp]
\centering
\caption{Step-4 IMU toe-off minus camera toe lift-off (build report, task 15).}
\label{tab:toeoff}
\begin{adjustbox}{max width=\textwidth}
{\small
\begin{tabular}{llrrrr}
\toprule
\textbf{rows} & \textbf{limb} & \textbf{n} & \textbf{IMU - camera, median (ms)} & \textbf{IQR (ms)} & \textbf{sensor height at the IMU toe-off, median (mm)} \\
\midrule
obstacle & lead & 292 & +67 & +17 to +167 & 62 \\
obstacle & trail & 300 & +233 & +150 to +350 & 356 \\
level & lead & 40 & $-$67 & $-$150 to $-$25 & 35 \\
level & trail & 42 & $-$25 & $-$133 to +17 & 31 \\
\bottomrule
\end{tabular}
}
\end{adjustbox}

\end{table}

\subsection{Files}

\begin{description}[leftmargin=0.4cm,style=nextline]
  \item[\vn{Crossing_Master_YoungAdults.csv}] 349 rows $\times$ 193 columns (303 obstacle crossings, 46 level walks).
  \item[\vn{Crossing_Curves_YoungAdults.csv}] Swing curves of both limbs, 51 points per signal.
  \item[\vn{Crossing_Aligned_YoungAdults.csv}] Both limbs' signals at 127 fixed samples around the trailing toe lift-off.
  \item[\vn{data_dictionary.md}, \vn{build_report.md}, \vn{exclusions_log.csv}] Definitions, checks and every cleaned or
  missing value with its reason.
\end{description}
